%============================================================ 8. 한계
\section{한계와 향후 연구}
\label{sec:limitations}

\subsection{시뮬레이션 한정}

본 연구의 모든 결과는 시뮬레이션에서 얻은 것이며 실해역 검증을 포함하지 않는다.
\S\ref{sec:deploy}에서 밝힌 성질들 --- 경유점과 동역학의 분리, 격자 상대 좌표 사용, 배치
통계 정규화의 부재 --- 은 실해역 이식에 유리한 조건이지만, 그 자체가 전이 성공의 증거는
아니다. 특히 다음 요소가 시뮬레이션에서 단순화되어 있다.

\begin{itemize}[leftmargin=*, itemsep=2pt, topsep=3pt]
\item \textbf{그물의 유체 거동과 전개 형상.} 본 모형에서 그물벽은 선분으로 표현되며
      표류를 거리 감쇠로 근사한다. 실제 전개는 곡선 형상의 면적을 이루며, 침강·조류에 의한 변형도 발생한다. 이 차이는 특히
      \emph{유효 차단 폭}을 과소 또는 과대 평가할 수 있다.
\item \textbf{추진기 얽힘 판정.} 적선이 그물에 접촉하면 무력화되는 것으로 처리하나, 실제
      얽힘 확률은 접근 속도와 각도에 의존한다.
\item \textbf{관측 완전성.} 적선의 위치와 진행 방향이 오차 없이 주어진다고 가정한다. 실제
      해상 감시에서는 탐지 누락, 위치 오차, 식별 지연이 발생한다.
\end{itemize}

특히 관측 불확실성은 점수맵 정책보다 지휘관 계층에 더 큰 영향을 줄 것으로 예상된다. 배정
판단은 무리 구조에 의존하는데, 무리 구성이 관측 오차로 흔들리면 \S\ref{sec:res_mechanism}에서
문제로 지적한 churn이 직접 증가하기 때문이다. 관측 불확실성 하에서의 계획 안정성은 향후
연구 과제이다.

\subsection{지연 측정의 교란 요인}

\S\ref{sec:res_latency}의 지연은 단일 GPU(RTX 4070, 12\,GB)에서 시뮬레이터·정책 추론과
언어모델 서빙을 함께 올린 구성에서 측정하였다. Qwen2.5 14B는 Q4 양자화 가중치만 약 9\,GB를
차지하므로 이 구성에서는 메모리 압박이 발생하며, 최댓값 296초는 가중치 오프로드를 시사한다.
따라서 14B에서 관측된 4.2\,\%의 주기 초과가 \emph{로컬 중형 모델 서빙의 본질적 성질인지,
이 자원 구성의 결과인지 본 실험은 구분하지 못한다}. 전용 가속기를 분리하거나 더 큰 메모리를
갖춘 환경에서 같은 측정을 반복하는 것이 필요하다.

다만 이 한계가 \S\ref{sec:arch}의 설계 논거를 무효화하지는 않는다. 온보드 배치에서 자원을
공유해야 하는 상황 자체가 현실적 제약이며, 비동기 구조는 그런 제약 아래에서 지연 꼬리가
제어 루프를 막는 것을 방지한다.

\subsection{폐쇄형 모델에 대한 의존과 재현성}

가장 좋은 성능을 보인 지휘관은 상용 API로 제공되는 폐쇄형 모델이다. 이런 모델은 버전이
예고 없이 갱신되거나 서비스가 종료될 수 있어, 그 수치의 재현이 장기적으로 보장되지 않는다.

본 연구는 이 위험을 구조적으로 완화하였다. 논문의 첫 번째 기여인 기동 계층의 성능
(\S\ref{sec:res_maneuver}, \texttt{heur\_unet} 조건)은 \textbf{언어모델을 전혀 사용하지
않으므로} 외부 API와 무관하게 재현 가능하다. 폐쇄형 모델은 지휘관 능력축의 상단 참조점으로만
기능한다. 또한 모든 계획 호출의 계측값(지연·폴백·커버리지·churn·교차·근거 길이)을 보존하였다. 다만 근거 원문은 저장하지 않아 사후 재검토가 제한된다.

그럼에도 \S\ref{sec:res_threshold}의 임계 위치는 시점 의존적인 수치이며, 절대적 기준으로
인용되어서는 안 된다.

\subsection{비교하지 않은 대안}

본 논문의 기동 계층 비교는 \emph{같은 행동공간 위에서} 이루어진다 --- 두 방법이 같은 격자,
같은 유효 마스크, 같은 거리 제약을 통과한다(\S\ref{sec:heuristic}). 이 통제가
\S\ref{sec:res_maneuver}의 비교를 성립시키는 조건이며, 그래서 기준 방법으로 기하
휴리스틱을 골랐다.

\textbf{표준 다중 에이전트 강화학습(MADDPG, QMIX, MAPPO)과는 비교하지 않았다.} 이유는
성능에 대한 자신감이 아니라 통제의 문제다. 이 계열을 관례대로 쓰면 행동이 연속 좌표나
에이전트별 이산 행동으로 표현되므로, 관측되는 차이에 \emph{표현}(연속 회귀 대 픽셀 지목)과
\emph{알고리즘}(가치망 유무, 크레딧 배분 방식)이 뒤섞인다. 그러면 \S\ref{sec:res_maneuver}가
답하려는 질문 --- 결정 규칙을 바꾸면 무엇이 달라지는가 --- 에 답할 수 없다. 공정한 비교는
그 알고리즘들을 같은 픽셀 지목 행동공간으로 옮겨 학습시키는 작업이며, 이는 별도의 연구
분량이다. 본 논문의 결과는 \emph{제안 표현이 최선임}을 주장하지 않는다.

\textbf{구성요소 절제 실험도 체계적으로 수행하지 않았다.} 개별 설계 선택에는 측정 근거를
붙였으나 --- 도달성 상한은 상한이 없을 때 선택 간격 중앙값이 504\unit{m}로 그물 길이
450\unit{m}를 넘고 초과율이 66.7\,\%였다는 실측에서, 잉여 배정의 층 분리는 집중 포메이션에서
배정의 과반이 그 경로로 결정된다는 관찰에서, 퍼텐셜 성형은 포획이 평가 창의 46.8\,\%에서만
발생한다는 계측에서 나왔다 --- 이는 각 요소가 \emph{왜 필요한가}에 대한 근거이지 \emph{얼마나
기여하는가}에 대한 측정이 아니다. FiLM 조건화, 좌표 채널, 자기회귀 두 번째 선택의 기여를
분리하려면 각각을 뺀 재학습이 필요하며, 본 논문은 이를 수행하지 않았다.

\subsection{규모와 일반화}

본 실험은 적 10척 대 방어정 3척의 고정된 규모에서 수행되었다. 점수맵 표현은 원리적으로 적선
수에 불변이고 파라미터 공유로 방어정 수에도 불변이지만, 그 불변성이 성능 유지로 이어지는지는
별도 검증이 필요하다. 특히 방어정 수가 늘면 배정 조합이 급격히 증가하여 지휘관 계층의 부담이
커지며, 이때 능력 임계가 어떻게 이동하는지는 측정되지 않았다.

또한 세 가지 공격 포메이션은 사전에 정의된 규칙 기반 행동이며, 방어 체계에 적응하는 지능적
적을 포함하지 않는다. 적이 그물벽 배치 패턴을 학습하여 회피하는 상황에서의 성능은 본 연구의
범위 밖이다.

\subsection{천장 지표}

\S\ref{sec:res_ceiling}에서 보였듯 집중 포메이션의 포획률은 기준선에서 이미 0.997이다. 이
구간을 포함한 통합 수치는 개선을 과소평가한다. 향후 평가에서는 난이도를 조절하여 천장을
피하거나(예: 적 수 증가, 방어정 수 감소), 천장 구간을 통합에서 분리 보고하는 것이 바람직하다.
